% Compatibilidad tipográfica de los desarrollos ampliados.
% Conserva los enunciados; utiliza los contadores del tratado integral.
\theoremstyle{definition}
\makeatletter
\@ifundefined{example}{\newtheorem{example}[theorem]{Ejemplo}}{}
\@ifundefined{axiom}{\newtheorem{axiom}[theorem]{Axioma}}{}
\makeatother
\newenvironment{teorema}{\begin{theorem}}{\end{theorem}}
\newenvironment{lema}{\begin{lemma}}{\end{lemma}}
\newenvironment{proposicion}{\begin{proposition}}{\end{proposition}}
\newenvironment{corolario}{\begin{corollary}}{\end{corollary}}
\newenvironment{definicion}{\begin{definition}}{\end{definition}}
\newenvironment{ejemplo}{\begin{example}}{\end{example}}
\newenvironment{observacion}{\begin{remark}}{\end{remark}}
\newenvironment{teoremarector}{\begin{theorem}}{\end{theorem}}
\newenvironment{tablacompleta}[1][tbp]
 {\begin{table}[#1]\centering\begin{minipage}{\linewidth}\centering}
 {\end{minipage}\end{table}}
\newenvironment{bloqueindivisible}
 {\par\addvspace{6pt}\noindent\begin{minipage}{\linewidth}}
 {\end{minipage}\par\addvspace{6pt}}
\providecommand{\notatabla}[1]{\par\vspace{5pt}{\footnotesize\raggedright #1\par}}
\providecommand{\NN}{\mathbb N}
\providecommand{\QQ}{\mathbb Q}
\providecommand{\CC}{\mathbb C}
\providecommand{\FF}{\mathbb F}
\providecommand{\Id}{\operatorname{id}}
\providecommand{\dd}{\mathrm d}
\providecommand{\Span}{\operatorname{span}}
\providecommand{\Cyl}{\operatorname{Cyl}}
\providecommand{\HMT}{\mathrm{HMT}}
\providecommand{\Tr}{\operatorname{Tr}}
\providecommand{\Spec}{\operatorname{Spec}}
\providecommand{\tr}{\operatorname{tr}}
\providecommand{\R}{\mathbb R}
\providecommand{\Z}{\mathbb Z}
\providecommand{\I}{\mathrm I}
\providecommand{\TRIT}{\{-1,0,+1\}}
\providecommand{\symbf}[1]{\boldsymbol{#1}}
\providecommand{\symbb}[1]{\mathbb{#1}}
\providecommand{\symcal}[1]{\mathcal{#1}}
\providecommand{\symup}[1]{\mathrm{#1}}
\newenvironment{HMTCompactSection}{\begingroup}{\par\endgroup}
\newenvironment{apunte}[1]{\par\Needspace{6\baselineskip}\begin{quote}\small\raggedright\noindent\textbf{Observación: #1.} }{\end{quote}}
\providecommand{\HMTNarrativePaper}{\clearpage\pagecolor{white}}
\providecommand{\HMTMathematicalPaper}{\clearpage\pagecolor{white}}
\providecommand{\HMTNarrativeHeading}[1]{\section*{#1}}
\providecommand{\HMTDivision}[4][]{\part{#3}\noindent #4\par}
\providecommand{\CatalogRouteReference}{Las secciones~\ref{vi:sec:vi-324-rutas} y~\ref{vi:sec:vi-inventario-metrologico} reúnen las 324 rutas orientadas y los 855 registros metrológicos de contraste.}
\DeclareUnicodeCharacter{03BD}{\ensuremath{\nu}}
\DeclareUnicodeCharacter{03B1}{\ensuremath{\alpha}}
\DeclareUnicodeCharacter{03B2}{\ensuremath{\beta}}
\DeclareUnicodeCharacter{03B4}{\ensuremath{\delta}}
\DeclareUnicodeCharacter{03B5}{\ensuremath{\varepsilon}}
\DeclareUnicodeCharacter{03B6}{\ensuremath{\zeta}}
\DeclareUnicodeCharacter{03B8}{\ensuremath{\theta}}
\DeclareUnicodeCharacter{03B9}{\ensuremath{\iota}}
\DeclareUnicodeCharacter{03BA}{\ensuremath{\kappa}}
\DeclareUnicodeCharacter{03BB}{\ensuremath{\lambda}}
\DeclareUnicodeCharacter{03BE}{\ensuremath{\xi}}
\DeclareUnicodeCharacter{03C5}{\ensuremath{\upsilon}}
\DeclareUnicodeCharacter{03C7}{\ensuremath{\chi}}
\DeclareUnicodeCharacter{03C9}{\ensuremath{\omega}}
\DeclareUnicodeCharacter{0393}{\ensuremath{\Gamma}}
\DeclareUnicodeCharacter{0394}{\ensuremath{\Delta}}
\DeclareUnicodeCharacter{0398}{\ensuremath{\Theta}}
\DeclareUnicodeCharacter{039B}{\ensuremath{\Lambda}}
\DeclareUnicodeCharacter{039E}{\ensuremath{\Xi}}
\DeclareUnicodeCharacter{03A0}{\ensuremath{\Pi}}
\DeclareUnicodeCharacter{03A3}{\ensuremath{\Sigma}}
\DeclareUnicodeCharacter{03A6}{\ensuremath{\Phi}}
\DeclareUnicodeCharacter{03A8}{\ensuremath{\Psi}}
\newcommand{\HMTresultado}[3]{%
  \par\addvspace{12pt}\Needspace{7\baselineskip}%
  \refstepcounter{theorem}\label{viii:#3}%
  \noindent{\normalfont\bfseries #1~\thetheorem. #2}\par\nobreak\vspace{5pt}}
\definecolor{ink}{HTML}{183D4A}
\definecolor{accent}{HTML}{8C3D2E}
\crefname{teorema}{teorema}{teoremas}
\crefname{lema}{lema}{lemas}
\crefname{proposicion}{proposición}{proposiciones}
\crefname{definicion}{definición}{definiciones}
\crefname{ejemplo}{ejemplo}{ejemplos}
\crefname{observacion}{observación}{observaciones}
% Fondo íntegramente blanco; jerarquía mediante tamaño, peso y espaciado.
\pagecolor{white}
\colorlet{hmtpaper}{white}
\colorlet{hmtlight}{white}
\colorlet{hmtlightblue}{white}
\colorlet{hmtlightgold}{white}
\colorlet{hmtlightred}{white}
\setcounter{tocdepth}{2}
\renewcommand{\cftsecfont}{\normalfont}
\renewcommand{\cftsecpagefont}{\normalfont}
\renewcommand{\cftsubsecfont}{\small\itshape}
\renewcommand{\cftsubsecpagefont}{\small}
